% 训练后仍未定的标签：被打散的六点集与阈值类对比（第5章5.3.2）。教学设定：n=3，|C|=2n=6，
% 输入 x_1<...<x_6 各有概率 1/6。训练样本见到 x_1、x_3、x_5，标签依次为 0、0、1（黄底标出）。
% (a) 类打散 C：未见点 x_2、x_4、x_6 的 8 种补全都可实现，三点全部未定，未定概率质量 3/6=1/2；
%     目标标签取独立公平硬币时，任何输出在每个未见点上的条件错误概率都是 1/2，平均风险至少 1/2*1/2=1/4。
% (b) 右侧判正的阈值 h_t(x)=I{x>=t}：x_3 为 0 推出 x_2 为 0，x_5 为 1 推出 x_6 为 1（浅色），只剩 x_4 未定，质量 1/6。
% 屏幕坐标：点 x_i 位于横坐标 22+11i（毫米）。
\begingroup
\begin{tikzpicture}[figure picture,x=1mm,y=1mm,
  one/.style={circle,fill=FigureBlue,draw=FigureBlue,inner sep=0pt,minimum size=2.8mm},
  zero/.style={circle,fill=white,draw=FigureRed,line width=.75pt,inner sep=0pt,minimum size=2.8mm},
  fone/.style={circle,fill=FigureBlue!40!white,draw=FigureBlue,dashed,line width=.6pt,inner sep=0pt,minimum size=2.8mm},
  fzero/.style={circle,fill=white,draw=FigureRed,dashed,line width=.6pt,inner sep=0pt,minimum size=2.8mm},
  free/.style={circle,fill=white,draw=FigureMuted,dashed,line width=.6pt,inner sep=0pt,minimum size=3.6mm,font=\figurefont\scriptsize\bfseries,text=FigureMuted},
  seen/.style={fill=FigureYellowFill,draw=FigureYellow,line width=.6pt,rounded corners=1mm}]
\node[figure heading,anchor=west] at (0,44) {假设空间};
\node[figure heading] at (61.5,44) {六个等概率输入（$n=3$）};
\node[figure heading] at (100,44) {未定质量};
\foreach \i in {1,...,6}{\node[figure note] at (22+11*\i,39) {$x_{\i}$};}
% 两行共同的已见位置
\foreach \i in {1,3,5}{\draw[seen] (22+11*\i-3.6,31.4) rectangle (22+11*\i+3.6,35.6);
                       \draw[seen] (22+11*\i-3.6,17.4) rectangle (22+11*\i+3.6,21.6);}
% (a) 被打散的集合
\node[anchor=west] at (0,33.5) {(a) 打散这六点};
\draw[FigureStroke] (30,33.5)--(93,33.5);
\node[zero] at (33,33.5) {}; \node[zero] at (55,33.5) {}; \node[one] at (77,33.5) {};
\node[free] at (44,33.5) {?}; \node[free] at (66,33.5) {?}; \node[free] at (88,33.5) {?};
\node at (100,33.5) {$3/6$};
% (b) 阈值
\node[anchor=west] at (0,19.5) {(b) 阈值类};
\draw[FigureStroke] (30,19.5)--(93,19.5);
\node[zero] at (33,19.5) {}; \node[zero] at (55,19.5) {}; \node[one] at (77,19.5) {};
\node[fzero] at (44,19.5) {}; \node[free] at (66,19.5) {?}; \node[fone] at (88,19.5) {};
\node at (100,19.5) {$1/6$};
\draw[FigureMuted,-{Stealth[length=1.6mm]},line width=.5pt] (52.5,16.3) to[bend left=35] (46.5,16.3);
\draw[FigureMuted,-{Stealth[length=1.6mm]},line width=.5pt] (79.5,16.3) to[bend right=35] (85.5,16.3);
% 图例
\node[figure note,anchor=west,align=left] at (0,6) {黄底：训练中出现的输入\quad 实心蓝：标签1\quad 空心红：标签0\\
虚线浅色：由已见标签推出\quad 虚线问号：仍可任意二选一};
\end{tikzpicture}
\endgroup
